%%% Methods overview (main text)

IMPG (Implicit Pangenome Graph) operates directly on pairwise alignments as implicit graph structures (Figure~\ref{fig:concept}A), providing three core capabilities: (a) transitive coordinate projection across all genomes in a pangenome, (b) alignment-based partitioning for distributed graph construction, and (c) on-demand graph materialization for regions of interest. Rather than materializing the full implicit graph, IMPG indexes alignment files using cache-oblivious interval trees \cite{dcjones2020} and traverses them at query time to discover both direct and transitive relationships. Given coordinates on any genome, IMPG projects them onto all other genomes, functioning as a pangenome-wide coordinate liftover. IMPG supports alignments as PAF files (uncompressed and BGZF-compressed), .1aln \cite{myers2025fastga}, and TPA \cite{tpa}---the latter two based on compact tracepoint representations of the alignments---and sequences as FASTA (uncompressed and BGZF-compressed) or AGC archives \cite{deorowicz2023agc}. This flexibility allows integration into existing workflows; supporting tracepoint-based formats and compressed genome archives enables population-scale queries on a single workstation. \textcolor{red}{On the HPRCv2 pangenome (466 haplotypes, 389\,M pairwise alignments in 466 per-file TPA files), IMPG indexes all alignments in X\,minutes and queries a 6\,Mbp region in under 2 seconds, with peak memory of X\,MiB (Supplementary Table~1). Transitive queries from a non-reference haplotype reach all 466 assemblies through multi-hop traversal---from X haplotypes (direct) to 466 (transitive)---demonstrating that the implicit index captures pangenome-wide connectivity (Extended Data Figure~\ref{fig:transitive_reach}). The complete ready-to-query representation of 466 haplotypes---per-file TPA alignments (23.3\,GiB), bidirectional indexes (14.7\,GiB), and AGC archive (3.1\,GiB)---totals 42.5\,GiB (Supplementary Table~1).} The entire human pangenome fits on a laptop disk and can be queried in seconds with minimal memory; no decompression or graph loading step is required.

As pangenome datasets grow to hundreds or thousands of haplotypes, monolithic graph construction becomes intractable: memory, computation, and time scale super-linearly with the number of sequences. Partitioning is essential, but existing approaches operate at whole-sequence granularity. Chromosome-based partitioning assigns entire contigs by chromosome identity; community-based approaches such as PGGB's Leiden clustering \cite{Garrison2024} assign entire contigs to communities based on sequence-similarity networks. Neither can split a contig: if one region of a contig is homologous to chromosome 1 and another to chromosome 5, the entire contig is forced into one partition.

IMPG introduces alignment-based partitioning, which operates at aligned-interval granularity (Figure~\ref{fig:concept}B). The \texttt{partition} command uses sliding windows of transitive queries to divide the pangenome into segments sharing alignment connectivity, and critically, it can split a single contig across multiple partitions at alignment boundaries. This naturally handles cases where different regions of a contig belong to different genomic contexts---including translocations, inter-chromosomal rearrangements, and fragmented assemblies---without requiring special-casing or forcing the entire contig into one partition.

This partitioning splits the pangenome into manageable chunks for any downstream analysis---similarity computation, sequence extraction, genotyping, or graph construction. For graph construction, each partition is processed independently by any tool---PGGB, vg \cite{Garrison2018}, or custom pipelines---distributing computation and memory across nodes; the \texttt{lace} command then merges the resulting subgraphs into a unified GFA or VCF. When a researcher needs to study region-specific variation---visualizing a complex structural variant, inspecting bubble structures, or calling variants against local topology---IMPG can materialize graphs on demand via three built-in engines: seqwish-based graph induction \cite{sweepga, garrison2023seqwish}, POA via SPOA \cite{Vaser2017}, and a recursive divide-and-conquer engine that chunks homologous sequences, builds POA sub-graphs in parallel, and laces them into a single variation graph. The \texttt{align} command supports this by generating sparsified alignment sets guided by Erd\H{o}s--R\'enyi connectivity thresholds \cite{erdos1959random}.

IMPG is already production infrastructure for structural genotyping. The COSIGT pipeline \cite{bolognini2026cosigt} uses \texttt{impg query} to project locus coordinates onto all haplotypes and \texttt{impg refine} to optimize locus boundaries, then constructs local pangenome graphs for genotyping, demonstrating that alignment-space operations can underpin population-scale genotyping of complex variation.
